\begin{ouverture}[calcul]{9}{Développer, réduire et démontrer}{Partie 2 · Algèbre et grandeurs}
{Distributivité simple · Facteur commun · Preuve · Contre-exemple}
\objectif{Développer et factoriser une expression simple.}
\objectif{Réduire une somme de termes de même nature.}
\objectif{Expliquer une propriété générale avec une lettre.}
\prerequis{Formules et priorités.}
\end{ouverture}
\pagestyle{otMain}\essentiel\fontsize{11}{15}\selectfont
\phantomsection\label{cours-9-1}
\essSection{1}{Distribuer ou regrouper}
\begin{propriete}Pour des nombres $k,a,b$, $k(a+b)=ka+kb$ et $k(a-b)=ka-kb$. Développer transforme un produit en somme ; factoriser transforme une somme en produit grâce à un facteur commun.\end{propriete}
\label{fig-cours-9-1-cours-distrib}\figureGeo{cours-distrib}

\begin{exemple}$4(x+3)=4x+12$. En sens inverse, $6x+18=6(x+3)$. Les deux écritures d’une même expression donnent les mêmes valeurs pour tous les $x$.\end{exemple}
\begin{attention}Dans $3(x+5)$, multiplier les deux termes par 3 : $3x+15$. L’écriture $3x+5$ oublie de multiplier le second terme.\end{attention}
\clearpage
\phantomsection\label{cours-9-2}
\essSection{2}{Réduire et raisonner}
\begin{propriete}$3x+5x=8x$ car on réunit huit fois le même nombre. $x+x=2x$, $1x=x$. Mais $3x+5$ ne se réduit pas à $8x$ : les termes n’ont pas la même nature.\end{propriete}
\begin{exemple}$2{,}5x+4+x+1=3{,}5x+5$. On regroupe les multiples de $x$ et les nombres sans lettre.\end{exemple}
\begin{savoirfaire}[Prouver pour toutes les valeurs]Le double de la somme de deux nombres consécutifs $n$ et $n+1$ vaut $2(n+n+1)=4n+2$. L’expression obtenue vaut pour tout entier $n$ : la lettre permet une démonstration générale.\end{savoirfaire}
\begin{definition}Un contre-exemple est un cas qui contredit une affirmation générale. Pour réfuter « $x+3=3x$ pour tout $x$ », choisir $x=1$ : 4 n’est pas égal à 3. Tester beaucoup de cas vrais ne remplace pas une preuve.\end{definition}

\clearpage\rubrique{Découverte}
\begin{blocLivre}{ex-9-1}\begin{exercice}[1]{Développer une somme}
Développer $3(x+4)$.
\end{exercice}\end{blocLivre}
\begin{blocLivre}{ex-9-2}\begin{exercice}[1]{Développer une différence}
Développer $5(x-2)$.
\end{exercice}\end{blocLivre}
\begin{blocLivre}{ex-9-3}\begin{exercice}[1]{Réduire des termes semblables}
Réduire $7x+2x+4$.
\end{exercice}\end{blocLivre}
\begin{blocLivre}{ex-9-4}\begin{exercice}[1]{Un facteur commun}
Factoriser $4x+20$.
\end{exercice}\end{blocLivre}
\begin{blocLivre}{ex-9-5}\begin{exercice}[1]{Avec des décimaux}
Réduire $1{,}5x+2x+3$.
\end{exercice}\end{blocLivre}

\rubrique{Entraînement}
\begin{blocLivre}{ex-9-6}\begin{exercice}[2]{Développer puis réduire}
Développer et réduire $2(x+5)+3x$.
\end{exercice}\end{blocLivre}
\begin{blocLivre}{ex-9-7}\begin{exercice}[2]{Une différence de termes}
Réduire $8x+7-3x+2$.
\end{exercice}\end{blocLivre}
\begin{blocLivre}{ex-9-8}\begin{exercice}[2]{Factoriser une différence}
Factoriser $7x-21$.
\end{exercice}\end{blocLivre}
\begin{blocLivre}{ex-9-9}\begin{exercice}[2]{Comparer deux programmes}
A calcule $4(x+2)$. B calcule $4x+8$. Ont-ils toujours le même résultat ?
\end{exercice}\end{blocLivre}
\begin{blocLivre}{ex-9-10}\begin{exercice}[2]{Une erreur à expliquer}
Aya écrit $2x+3=5x$. Donner un contre-exemple et corriger.
\end{exercice}\end{blocLivre}

\rubrique{Problèmes}
\begin{blocLivre}{ex-9-11}\begin{exercice}[3]{Un périmètre général}
Un rectangle a pour côtés $x+4$ et $x+1$ cm. Exprimer son périmètre sous forme développée et réduite.
\label{fig-ex-9-11}\figureGeo{dist-11}
\end{exercice}\end{blocLivre}
\begin{blocLivre}{ex-9-12}\begin{exercice}[3]{Une aire découpée}
Un rectangle de hauteur 3 cm a deux parties de largeurs $x$ et 5 cm. Donner deux expressions de son aire et expliquer leur égalité.
\label{fig-ex-9-12}\figureGeo{dist-12}
\end{exercice}\end{blocLivre}
\begin{blocLivre}{ex-9-13}\begin{exercice}[3]{Trois entiers consécutifs}
On considère $n$, $n+1$ et $n+2$. Montrer que leur somme est trois fois celui du milieu.
\end{exercice}\end{blocLivre}
\begin{blocLivre}{ex-9-14}\begin{exercice}[3]{Un programme invariant}
Choisir un nombre $x$, lui ajouter 6, doubler le résultat puis retirer le double du nombre choisi. Prouver que le résultat est toujours 12.
\end{exercice}\end{blocLivre}
\begin{blocLivre}{ex-9-15}\begin{exercice}[3]{Une conjecture trompeuse}
Un tableur contient les entiers de 1 à 10 dans la colonne A. Écrire en B2 la formule pour calculer $x^2+x+11$ à partir de A2, puis la recopier. Les premières valeurs suggèrent-elles que le résultat est toujours premier ? Examiner particulièrement $x=1$, 2 et 10.
\end{exercice}\end{blocLivre}

\clearpage
